\documentclass[11pt,compress,t,notes=noshow, xcolor=table]{beamer}

\input{../../style/preamble}
\input{../../latex-math/basic-math}
\input{../../latex-math/basic-ml}

\newcommand{\xleft}{x_{\text{left}}}
\newcommand{\xright}{x_{\text{right}}}
\newcommand{\xnew}{x_{\text{new}}}
\newcommand{\xbest}{x_{\text{best}}}
\newcommand{\fleft}{f_{\text{left}}}
\newcommand{\fright}{f_{\text{right}}}
\newcommand{\fnew}{f_{\text{new}}}
\newcommand{\fbest}{f_{\text{best}}}
\newcommand{\xii}{x^{(i)}}

\title{Optimization in Machine Learning}

\begin{document}

\titlemeta{
Univariate optimization
}{
Brent's method
}{
figure_man/quadratic-title.png
}{
\item Quadratic interpolation
\item Brent's procedure
\vspace{0.5cm}
\item[] \hspace{-0.7cm}Slides based on \furtherreading{PRESS92} (Ch. 10.2)
}

\begin{frame}{Quadratic Interpolation}
Similar to golden section search but select $\xnew$ differently: \\
minimum of parabola through $(\xleft, \fleft), (\xbest, \fbest), (\xright, \fright)$
\foreach \i in {1, 2, 3, 4, 5}{
\only<\i>{%
  \begin{center}
    \includegraphics[width=\textwidth]{figure/quadratic-\i.pdf}\\[0.3em]
    {\footnotesize Left: Fit parabola (dashed) and propose its minimum (red) as new point.\\
    Middle: Switch $\xnew$ and $\xbest$ if $f(\xnew) < \fbest$. Right: New interval.}
  \end{center}%
}}
\end{frame}


\begin{framei}[sep=L]{Quadratic interpolation comments}
\item Much faster than GSS for smooth $f$: well-approximated by a parabola near the minimum $\Rightarrow$ QI jumps directly toward the minimum
\item But: not as robust as GSS, the following may happen:
\comment[NOTE]{BB}{hier wäre es päter mal gut diese fälle einfach mal zu zeigen (interpolation)}
\comment[TODO]{TH}{not sure whether case 2 and 3 are even possible under our assumptions?}

\begin{itemizeL}
\item Algorithm suggests very similar $\xnew$ in each step
\item $\xnew$ outside of search interval
\item Parabola degenerates to line and no real minimum exists
\end{itemizeL}
\item Algorithm must then abort, finding the minimum is not guaranteed
\end{framei}


\begin{framei}{Brent's method}
\item Brent's algorithm \furtherreading{BRENT1973} alternates between golden section search and quadratic interpolation as follows:
\begin{itemizeM}
\item Quadratic interpolation step acceptable if:\\
(i) $\xnew$ falls within $[\xleft, \xright]$\\
(ii) $\xnew$ sufficiently far away from  $\xbest$ \\
{\footnotesize (Heuristic: Less than half of the movement of step before last)}
\item Otherwise: proposal via GSS
\end{itemizeM}
\item GSS stabilizes unstable steps and guarantees \textbf{at least linear convergence} for any unimodal $f$
\item Faster convergence (quadratic interpolation) on well-behaved (smooth) $f$: \textbf{superlinear convergence} of order $\approx 1.325$
\item Note: Brent's algorithm also exists with other robust methods than GSS, e.g., bisection method
\end{framei}


\begin{framei}{Example: MLE Poisson}
\item Poisson PMF: $p(k|\lambda) := \P(x = k) = \frac{\lambda^k \cdot \exp(-\lambda)}{k!}$
\item Negative log-likelihood for $n$ samples:
\begin{equation*}
- \ell(\lambda, \mathcal{D}) = - \log \prodin  p(\xii | \lambda) =  - \sumin \log p(\xii | \lambda)
\end{equation*}
\imageC[0.7]{figure/poisson.pdf}
GSS and Brent converge to global min at $x^\ast \approx 1$ \\
But GSS needs $\approx 45$ iters, Brent only $\approx 15$ for same tolerance
\end{framei}

\endlecture
\end{document}
